\section{Experimental design}
\label{sec:experiments}
\subsection{Development protocol and landmark evaluation}
The 801 development patients were partitioned into five patient-level folds. Fold 1 used 640 training and 161 validation patients; folds 2--5 each used 641 training and 160 validation patients. The same 175-patient holdout was evaluated by every fold model and every variant. Thus the five test evaluations per variant are repeated evaluations of overlapping trained models on the same patients, rather than five independent test cohorts. Patient identities were disjoint between training, validation, and holdout within the current split.

At each training epoch, a patient contributed one randomly selected available T0, T1, or T2 prefix. Only MRI observations available at or before the selected landmark were classifier inputs. Later observed visits, including T3, could supply latent targets for the auxiliary loss on training patients. Missing targets were masked. The 23 patients with only T0 remained eligible for pCR training but supplied no future-state loss. At evaluation, all three landmarks were evaluated for the 175 holdout patients, each of whom had T0, T1, and T2. T3 was not an evaluation landmark and was not used as observed input for these predictions.

\subsection{Model comparisons}
The principal comparison comprises three variants trained with seed 44 and the same five development folds. Each variant is evaluated as an equal-weight ensemble of its five fold-specific models. \nojepa{} uses direct pCR readout and a zero future-state loss weight. \direct{} uses the same direct readout with a future-state loss weight of 0.5. \fullmodel{} retains that auxiliary weight and replaces the direct readout with a concatenation of direct pCR features, attention-pooled predicted terminal state, and attention-pooled predicted change. All variants retain the shared image encoder, clinical inputs, assigned-arm encoding, fusion architecture, and training schedule.

The \direct{}--\nojepa{} contrast most directly tests adding future-state supervision. The \fullmodel{}--\direct{} contrast jointly changes the use of forecast features, spatial pooling, and classifier input width. It cannot isolate attention alone. For that purpose, the supplementary \meanmodel{} comparison uses forecast readout with mean spatial pooling and otherwise matches the forecast-attention configuration. All manuscript-facing pCR results use seed 44 and average the five fold probabilities separately for each patient and landmark before computing metrics. No VQ-latent reconstruction-head, change-weighted-loss, or state-conditioned-head variants are included.

\subsection{Implementation and optimization}
The network uses a model width of 768, eight layers in each of the spatial encoder, multimodal fusion module, and predictor, 16 attention heads with 64 dimensions per head, and feed-forward width 3,072. The 1,024-dimensional concatenated attention-head representation is projected back to the 768-dimensional residual stream. The spatial token grid is $2\times4\times4$. The configuration uses eight state queries, four pCR queries, dropout 0.3, and an EMA coefficient of 0.999. The pretrained convolutional encoder is fine-tuned jointly with the remaining network. Its training-patient provenance requires additional documentation, as discussed in Section~\ref{sec:limitations}.

We use AdamW \citep{adamw2019} with an initial base learning rate of $10^{-4}$ for both the encoder and the main network, weight decay 0.05, batch size 16, no gradient accumulation, and gradient-norm clipping at 1.0. Training uses BF16 mixed precision. A ten-epoch warm-up increases the learning rate from $10^{-5}$ to $10^{-4}$, followed by cosine decay to $10^{-6}$ at epoch 100. Training is capped at 200 epochs, with early stopping after 50 epochs without improvement. The best checkpoint is selected independently in each fold using validation patient-equal-prefix F1 at threshold 0.5. This is an F1 selection protocol, even though AUROC is a principal reporting measure. Deterministic execution was not enabled, so a fixed seed does not imply bitwise reproducibility across hardware or software changes.

\subsection{Metrics, aggregation, and patient-level uncertainty}
We report 14 complementary metrics: AUROC, non-interpolated average precision (AP), F1, accuracy, sensitivity, specificity, positive and negative predictive value (PPV/NPV), balanced accuracy, Matthews correlation coefficient (MCC), Youden index, negative log-likelihood (NLL), Brier score, and expected calibration error (ECE) with ten equal-width probability bins. Confusion-matrix counts are additionally retained. Thresholded primary metrics use 0.5. MCC and balanced accuracy help interpret performance under class imbalance~\citep{chicco2020mcc}; positive prevalence is $56/175=0.32$. No holdout threshold optimization or probability recalibration was performed.

Let $p_{if\ell}$ be the probability from the seed-44 model trained in fold $f$ for patient $i$ at landmark $\ell$. For each variant, the five fitted models form one probability ensemble,
\begin{equation}
\begin{aligned}
 \bar p_{i\ell}&=\frac{1}{5}\sum_{f=1}^{5}p_{if\ell},\\
 M_{\ell}&=M\bigl(\{y_i,\bar p_{i\ell}\}_{i=1}^{175}\bigr),\\
 M_{\mathrm{macro}}&=\frac{1}{3}\sum_{\ell=0}^{2}M_{\ell}.
\end{aligned}
\end{equation}
The probability average is taken before computing every metric, curve, confusion matrix, or subgroup summary. Classification applies threshold 0.5 to $\bar p_{i\ell}$. Macro is an equal mean of the three landmark metrics, including F1; it is not a metric recomputed after pooling landmark rows. This estimand describes a five-model predictor trained with seed 44, rather than the average performance of a single fold model. The same ensemble rule is applied to every variant.

For exploratory uncertainty estimates, we draw 2,000 bootstrap samples of the 175 patients with replacement. Every draw is shared across variants and landmarks. Within each resample, metrics are evaluated on the fixed ensemble probabilities, and 2.5th and 97.5th percentiles form the intervals. These intervals are not adjusted for multiple comparisons. They condition on the five fitted models per variant and do not capture retraining, seed selection, or model selection. Neither the five constituent models nor the three repeated landmarks are treated as independent patient samples.

ROC and precision--recall curves are computed directly from the ensemble probabilities. ROC bands use a common false-positive-rate grid; PR bands use a common recall grid. AUROC and non-interpolated AP are calculated directly from the scores; grid interpolation may change the area of a plotted curve. Reliability diagrams divide ensemble probabilities into five equal-count rank groups of 35 patients, distinct from the ten fixed bins used for ECE.

\subsection{Exploratory paired significance checks}
The main significance checks test a zero difference in the reported population metric, conditional on the fitted ensembles and evaluation policy. AUROC comparisons use a two-sided paired DeLong asymptotic normal test~\citep{delong1988}. For Macro AUROC, we first average the three landmark-specific positive-patient placement vectors and, separately, the three negative-patient placement vectors within each model. The variance of the paired model difference is then estimated from these averaged placements, retaining both model and landmark dependence. The same paired procedure compares T2 with T0 within a model. Neither folds nor repeated landmarks are independent patient samples.

For the other 13 metrics, we use approximate, null-centered paired patient-bootstrap tests~\citep{davison1997}, with $B=99{,}999$ resamples and random-number seed 20261103. Each sample draws 175 patients with replacement, sharing the patient draw across all models and landmarks. Metrics and their signed differences are recomputed on the fixed ensemble scores, using threshold 0.5 for thresholded metrics and averaging separately computed landmark metrics for Macro. Let $\Delta_b^*$ be a bootstrap difference and $\Delta_{\mathrm{obs}}$ its observed value. The two-sided Monte Carlo estimate is
\begin{equation}
 p=\frac{1+\sum_{b=1}^{B}\mathbb{I}\bigl(|\Delta_b^*-\Delta_{\mathrm{obs}}|\geq|\Delta_{\mathrm{obs}}|\bigr)}{B+1}.
\end{equation}
Ties count as at least as extreme. The plus-one Monte Carlo convention avoids reporting zero; it does not make this bootstrap test finite-sample exact. AP, F1, and ECE are non-smooth statistics, for which finite-sample bootstrap approximations warrant caution. No model is retrained.

The principal reporting family contains nine tests: the three pairwise comparisons among \nojepa{}, \direct{}, and \fullmodel{}, each for Macro AUROC, AP, and F1. The secondary family contains 81 tests: 27 landmark-specific AUROC/AP/F1 comparisons among the principal models; 33 comparisons of the other 11 Macro metrics; 12 \fullmodel{}--\meanmodel{} comparisons for AUROC/AP/F1 at T0, T1, T2, and Macro; and nine T2--T0 AUROC/AP/F1 comparisons within the three principal models. Holm adjustment~\citep{holm1979} is applied separately within each family. The aggregate CSV additionally reports a sensitivity analysis adjusting all 90 tests together; separate family adjustments do not control the familywise error rate over their union. We use $p_{\mathrm{Holm}}<0.05$, with $*$, $**$, and $***$ denoting values below 0.05, 0.01, and 0.001, respectively, and ns otherwise.

The original score-swap randomization checks remain fully exported as a sensitivity analysis. They exchange whole T0--T2 score trajectories within patients, or T0/T2 scores for longitudinal contrasts, using 99,999 randomizations and plus-one Monte Carlo $p$ values~\citep{phipson2010}. Their null is conditional exchangeability of scores, which is stronger than equality of a metric; they do not determine the main tables' significance markers. After inspecting initial score-swap results, we revised the main tests to target metric equality. Neither this method choice nor the seed, ensemble policy, or reporting families was prospectively specified.

The existing 2,000-resample percentile intervals remain unadjusted and are not exact inversions of these tests, so their exclusion of zero need not agree with adjusted $p$ values. We do not infer $p$ values from the intervals. All significance checks remain exploratory, including any result passing adjustment; they exclude retraining and retrospective selection uncertainty and do not establish stability across training seeds. OOF-threshold, subgroup, decision-curve, and training-probe summaries receive no additional significance tests.

\subsection{Operating-point and hypothetical decision analyses}
As a threshold sensitivity analysis, each development patient's out-of-fold (OOF) probability is supplied only by the fold model for which that patient was in validation. These patients also informed their fold-specific best-checkpoint selection, so the OOF scores are not a nested or independent calibration set. The disjoint validation predictions are concatenated separately at T0, T1, and T2, and Youden index is maximized on the resulting OOF scores, with the highest threshold resolving ties. The threshold is then applied to holdout ensemble probabilities. The OOF score for a development patient comes from one model, whereas the holdout score averages five models; their score distributions can differ. This analysis therefore examines threshold transfer and does not supply independent validation or calibration of the five-model ensemble. No holdout outcome selects a threshold. Bootstrap intervals hold the OOF-selected threshold fixed.

We additionally compute exploratory decision curves~\citep{vickers2006dca}. At risk threshold $q$, the generic positive action is applied when predicted pCR probability is at least $q$. Its mathematical net benefit is
\begin{equation}
 \mathrm{NB}(q)=\frac{\mathrm{TP}(q)}{n}
 -\frac{\mathrm{FP}(q)}{n}\frac{q}{1-q}.
\end{equation}
We calculate this directly from the five-fold ensemble probabilities over thresholds 0.05--0.80 and average the three landmark net benefits for the macro summary. Treat-all and treat-none are references. The calculation uses uncalibrated probabilities and does not validate a particular escalation, de-escalation, or other management action.

\subsection{Exploratory subgroup and mechanism analyses}
We examine released HR/HER2 combinations, age below or at least 50 years, native source namespace, and saved predicted-component crop coverage. Age 50 is an exploratory reporting cutpoint. Subgroup AUROC/AP are suppressed when fewer than ten patients or fewer than five outcomes in either class are available. Paired within-subgroup patient resampling preserves model and landmark pairing of the ensemble predictions; single-class draws are omitted for AUROC intervals and valid-draw counts retained. Intervals are unadjusted, and no subgroup-interaction hypothesis is tested. Source namespaces share the same clinical release and are not independent external cohorts.

Mechanism analysis is restricted to archived fixed-training-probe diagnostics collected every five epochs. These probes sample up to 32 training patients and are not full validation or holdout forecast evaluations. Unlike the missing-target-masked training objective, the archived forecast diagnostic includes only continuous observed visit chains: the real landmark, target, and every intermediate canonical visit must be available. A gap excludes that landmark--target pair, and eligible-patient counts therefore vary by horizon. We report a descriptive common epoch 20 and common-epoch trajectories. Forecast skill is $1-\mathrm{L1}_{\mathrm{forecast}}/\mathrm{L1}_{\mathrm{copy}}$ within the same snapshot's EMA representation; the copy forecast holds the observed landmark state constant. This normalization avoids comparing raw latent errors between separately learned teacher coordinates. Predictor gradient norms from the pCR and weighted rollout objectives are descriptive checks of gradient routing, not independent statistical samples or clinical effects. These archived diagnostics are averaged across the five seed-44 fold models after computing each statistic within its own EMA coordinates. They do not ensemble latent states and should be distinguished from the probability-ensemble pCR evaluation.

\subsection{Scope of the selected-seed ensemble analysis}
The choice of seed 44 and probability ensembling was made retrospectively after multiple-seed holdout results were available. It was not a prospectively registered seed or aggregation policy. Applying the same seed and ensemble rule to every model preserves a matched comparison, but does not remove selection bias or characterize training randomness. A confirmatory study should lock the model, seed, ensembling, and operating-point policies before evaluating independent patients.
